% Abhakseite „Das kann ich“ – Zone nur im Gesamt (\mitzone)
%% TODO Ich-kann-Sätze: Platzhalter wie in den Aufgabendateien
\begin{abhakseite}
\ifdefined\mitzone
\abhakgruppe{Kennst du schon}
\abhak{1}{Punkte im ebenen Koordinatensystem mit vier Quadranten lesen und eintragen, Vorzeichen der Koordinaten deuten}
\abhak{2}{Schrägbilder von Quadern, Prismen und Pyramiden lesen und skizzieren, Netze aufklappen}
\abhak{3}{Verbindungsvektor bilden (Spitze minus Fuß), Vektoren addieren und vervielfachen, den Betrag berechnen}
\abhak{4}{Satz des Pythagoras anwenden und mit Wurzeln rechnen (Wurzel isolieren, dann quadrieren)}
\abhak{5}{Flächenformeln (Dreieck, Trapez), Prozentrechnung und Strahlensätze für die Sachanhänge}
\abhak{6}{Skalarprodukt in Koordinaten berechnen und Orthogonalität als „Skalarprodukt null“ erkennen}
\abhak{7}{Eigenschaften der Dreiecke und Vierecke kennen (gleichschenklig, gleichseitig, rechtwinklig; Parallelogramm, Rechteck, Raute, Quadrat, Trapez, Drachenviereck) und den Satz des Thales}
\abhak{8}{Kollinearität zweier Vektoren prüfen (Vielfaches)}
\abhak{9}{Verbindungsvektor bilden (Spitze minus Fuß), Vektoren addieren und vervielfachen, den Betrag berechnen – Fehler finden}
\abhak{10}{Verbindungsvektor bilden (Spitze minus Fuß), Vektoren addieren und vervielfachen, den Betrag berechnen – selbst rechnen}
\fi
\abhakgruppe{Einheit 1 · Punkte darstellen und Lage lesen}
\abhak{11}{Darstellen und Lage lesen}
\abhak{12}{Darstellen und Lage lesen – Fehler finden, Begründen}
\abhakgruppe{Einheit 2 · Streckenlängen, Mittelpunkt und Teilpunkte}
\abhak{13}{Streckenlänge, Mittelpunkt, Teilpunkte}
\abhak{14}{Streckenlänge, Mittelpunkt, Teilpunkte – Fehler finden, Begründen}
\abhakgruppe{Einheit 3 · Dreiecke nachweisen}
\abhak{15}{Dreiecke nachweisen}
\abhak{16}{Dreiecke nachweisen – Fehler finden, Begründen}
\abhakgruppe{Einheit 4 · Vierecke nachweisen}
\abhak{17}{Vierecke nachweisen}
\abhak{18}{Vierecke nachweisen – Fehler finden, Begründen}
\abhakgruppe{Einheit 5 · Körper im Koordinatensystem und Drehungen}
\abhak{19}{Körper und Drehungen}
\abhak{20}{Körper und Drehungen – Fehler finden, Begründen}
\end{abhakseite}
